\documentclass[11pt,a4paper]{article}
\usepackage[utf8]{inputenc}
\usepackage[T1]{fontenc}
\usepackage{amsmath,amssymb}
\usepackage[margin=2.7cm]{geometry}
\usepackage{hyperref}
\hypersetup{colorlinks=true,linkcolor=blue,citecolor=blue,urlcolor=blue}

\title{The radion constrains the dark dimension:\\
an exclusion map for Goldberger--Wise stabilization}

\author{Hicham Boufourou\\[2pt]
\small Independent researcher, Ingelmunster, Belgium\\
\small \texttt{hicham.boufourou@hotmail.com}}

\date{DRAFT --- 2026. \textbf{Conditional on the official digitization of the
$\alpha_{95}$ curves; not for submission before that input is fixed.}}

\begin{document}
\maketitle

\begin{abstract}
The dark dimension scenario predicts a single extra dimension of
micron scale, with the size fixed by the observed cosmological
constant. Its authors note that the associated radion must be
decoupled from matter to avoid fifth-force constraints, and describe
this as a question to be explored. We take the question at face value
for a background stabilized by a Goldberger--Wise bulk scalar, in
which the radion coupling and range are \emph{derived} from the
background rather than assumed. Confronting the combined signal ---
radion Yukawa plus the full Kaluza--Klein tower --- with sub-millimetre
torsion-balance data, we obtain an exclusion map in the space of
background parameters. In the region we scan, the micron-scale
benchmark is excluded by its own radion, and the surviving window is
nanometric in the extra-dimension radius while the radion range stays
in the tens of microns. We claim the map and its verdict, not the
coupling: $\alpha_r=1/3$ for one extra dimension is standard and is
cited, not derived here.
\end{abstract}

\section{Motivation and what is claimed}

The dark dimension scenario~\cite{MVV} places one extra dimension at
the micron scale, with $m_{\rm KK}\sim\Lambda^{1/4}$, and a
five-dimensional Planck scale around $10^9\,$GeV. Its phenomenology
has been developed in several directions, notably Kaluza--Klein
gravitons as dark matter~\cite{GMOV}. The radion, however, has received
much less attention: \cite{MVV} states that one must decouple the
radion from the matter sector to avoid trouble with fifth-force
constraints --- either by making it heavy or by suppressing its
couplings --- and calls this a question to be explored further. A recent
realization stabilizing the radion by the Casimir effect finds it too
light to survive Solar-System tests and appeals to a screening
mechanism~\cite{Cas2026}.

Two facts frame what follows. First, the radion Yukawa strength for
$n$ extra dimensions is standard, $\alpha=n/(n+2)$, hence
$\alpha_r=1/3$ for $n=1$; this is the value used by the torsion-balance
collaborations themselves~\cite{AHH}, and it follows from the radion of
a single extra dimension being a Brans--Dicke scalar with $\omega=0$.
Second, the Kaluza--Klein tower contributes $\alpha_{\rm KK}=8n/3$ with
range $R$~\cite{KS}. Neither is claimed here.

\medskip
\noindent\textbf{What is claimed.} For a background stabilized by a
Goldberger--Wise bulk scalar, we derive the radion range $\lambda_r$
and the correction to $\alpha_r$ \emph{from the background parameters},
and we produce the resulting exclusion map. The result --- that the
micron-scale benchmark is excluded by its own radion, and that the
surviving window is nanometric --- is the claim.

\section{Setup}

The background is the one whose scalar sector has been shown to be free
of ghosts and tachyons~\cite{PaperA}: a two-brane interval with a
$\mathbb{Z}_2$-symmetric Goldberger--Wise scalar,
$\sigma_0(y)=(q/m_\sigma)\sinh[m_\sigma(y-L/2)]/\cosh(x/2)$ with
$x=m_\sigma L$. The dimensionless parameters are $x$, the ratio
$m_\sigma/m_\Phi$, the stabilizer amplitude $q$, and
$\eta=2\Lambda_5/q^2$.

\paragraph{Radion coupling.} In the Einstein frame,
\begin{equation}
\alpha_r=\tfrac13\Big[1+c_\alpha(\beta L)^2\Big],\qquad
c_\alpha=-0.634,
\label{eq:alpha}
\end{equation}
where $\beta L$ measures the backreaction. The leading $1/3$ is
\cite{AHH}; the correction $c_\alpha$ is derived here in the
moduli-space approximation and is tagged as such.

\paragraph{Radion mass and range.}
$m_r=2\sqrt{\eta x\tanh(x/2)}\;q/\sqrt{12M_5^3}$ up to a computed
correction of order one percent from backreaction, and
$\lambda_r=1/m_r$.

\paragraph{The signal.} The deviation from the inverse-square law is
\begin{equation}
\boxed{\;\Delta(r)=\alpha_r\,e^{-r/\lambda_r}
+\frac{8}{3}\sum_{n\ge1}e^{-nr/R}\;}
\label{eq:signal}
\end{equation}
Both legs are present; neither may be dropped. In particular the
single-Yukawa reading $(\alpha,\lambda)=(8/3,R)$ is a visualization,
not the test.

\section{The exclusion map}

We scan the dimensionless background parameters, never $R$ or
$\lambda_r$ directly, and we never tune to pass a curve. Theoretical
admissibility requires the linearity of the stabilizing well to better
than one percent, backreaction $0.49(\beta L)^2<10^{-2}$,
$\beta L<\pi/8$, and effective-theory validity. The experimental
verdict compares \eqref{eq:signal} to the published response in its own
domain.

\paragraph{Result (conditional).} In the scanned region the surviving
window is
\begin{equation}
R^\ast\in[0.003,\,0.056]\ \mu\mathrm{m},\qquad
\lambda_r\in[6,\,20]\ \mu\mathrm{m},
\end{equation}
with an amber band extending to $\lambda_r\simeq70\,\mu$m and
$R^\ast\simeq0.23\,\mu$m. The historical micron benchmark
($R^\ast\simeq7.4\,\mu$m) fails by a factor of order $260$ --- and it
fails on the radion leg, not the tower.

\section{Declared conditionals}

\begin{enumerate}
\item \textbf{Blocking.} The verdict is conditional on the official
digitization of the $\alpha_{95}$ response curves. The screening curve
used to build the map is an approximation, anchored on one published
point and uncertain by factors of two to three elsewhere; it is
adequate to sort the parameter space, and it is \emph{not} adequate to
publish an exclusion. This paper is not to be submitted before that
input is replaced.
\item $c_\alpha$ is derived in the moduli-space approximation.
\item The nanometric window displaces the Kaluza--Klein tower from
meV to the eV range. Any dark-matter role for the tower must therefore
be rebuilt at that scale; nothing is inherited.
\item $\alpha_r=1/3$, $\alpha_{\rm KK}=8n/3$, and the micron-scale
scenario itself are cited, not claimed.
\end{enumerate}

\begin{thebibliography}{9}
\bibitem{MVV} M.~Montero, C.~Vafa and I.~Valenzuela,
\emph{The dark dimension and the Swampland}, JHEP \textbf{02} (2023) 022,
\href{https://arxiv.org/abs/2205.12293}{arXiv:2205.12293}.
\bibitem{GMOV} E.~Gonzalo, M.~Montero, G.~Obied and C.~Vafa,
\emph{Dark dimension gravitons as dark matter}, JHEP \textbf{11} (2023) 109,
\href{https://arxiv.org/abs/2209.09249}{arXiv:2209.09249}.
\bibitem{Cas2026} \emph{Casimir energy stabilization of Standard Model
landscape in the dark dimension}, JHEP \textbf{02} (2026) 156.
\bibitem{AHH} E.~G.~Adelberger, B.~R.~Heckel, C.~D.~Hoyle et al.,
\emph{Particle physics implications of a recent test of the
gravitational inverse square law},
\href{https://arxiv.org/abs/hep-ph/0611223}{arXiv:hep-ph/0611223}.
\bibitem{KS} A.~Kehagias and K.~Sfetsos,
\emph{Deviations from the $1/r^2$ Newton law due to extra dimensions},
Phys.\ Lett.\ B \textbf{472} (2000) 39,
\href{https://arxiv.org/abs/hep-ph/9905417}{arXiv:hep-ph/9905417}.
\bibitem{Lee} J.~G.~Lee et al., \emph{New test of the gravitational
$1/r^2$ law at separations down to 52 $\mu$m}, Phys.\ Rev.\ Lett.\
\textbf{124} (2020) 101101.
\bibitem{PaperA} H.~Boufourou, \emph{A regular energy identity for the
scalar sector of stabilized two-brane models} (companion paper).
\end{thebibliography}

\end{document}
